% GENERATED FILE. Edit canonical lessons, then run npm run generate:books.
% canonical-source-hash: fnv1a64:ffd5b0846866a81c
% canonical-lessons: RU-C387-znanie, RU-C387-umenie, RU-C387-usvaivat, RU-C387-otchislyat, RU-C387-obrazovannyy

\chapter{Knowledge and skills}
\label{ch:ru-knowledge-and-skills}
\hlchaptermodality{387}

\begin{chapteropening}
\textbf{By the end of this chapter:} \emph{I can talk about knowledge and skills, say how well someone takes things in, call someone educated, and say a student was expelled.}

Complete the last lesson of chapter 387: I can talk about knowledge and skills, say how well someone takes things in, call someone educated, and say a student was expelled.
\end{chapteropening}

\section[\texorpdfstring{\ru{знание} (znániye)}{знание (znániye)}]{\ru{знание} (znániye) — knowledge}

\label{lesson:RU-C387-znanie}



\begin{quote}
Before the new one: say the Russian for a concept, then the Russian for a term.
\end{quote}

\subsection*{You'll want to know: \ru{знание}}
\textbf{\ru{знание}} (\emph{znániye}) — \textquotedblleft{}knowledge\textquotedblright{}.

\begin{quote}
\ru{Для} \ru{этой} \ru{работы} \ru{нужны} \ru{знания} \ru{языков}.
\end{quote}

\begin{quote}
\emph{Dlya étoy rabóty nuzhný znániya yazykóv.}
\end{quote}

\begin{quote}
This job needs a knowledge of languages.
\end{quote}

Neuter; usually plural, \emph{\ru{знания}}: \emph{\ru{хорошие} \ru{знания}}.

\begin{cousinweb}
\textbf{\ru{знание}} is built on \emph{\ru{знать}} \textquotedblleft{}to know\textquotedblright{}, inherited through Old East Slavic from Proto-Slavic \emph{*znati} (EtymDB).
\end{cousinweb}

\begin{grammarlens}[title={What you've built}]
One more word for telling, explaining and describing at B1.
\end{grammarlens}

\subsection*{Your turn}
\begin{itemize}
\raggedright
  \item \emph{Say it:} \emph{znániye}
  \item \emph{Say it:} \emph{znániye}, once more
  \item \emph{Say it:} a sentence of your own with \emph{znániye}
  \item \emph{Recall:} say the Russian for a concept, then the Russian for a term, then say \emph{znániye} again
\end{itemize}

\begin{tcolorbox}[breakable,colback=teal!4,colframe=teal!35!black,title={Before you move on}]
What does \ru{знание} mean? (\textquotedblleft{}Knowledge\textquotedblright{}.) Say it once more.
\end{tcolorbox}
\section[\texorpdfstring{\ru{умение} (uméniye)}{умение (uméniye)}]{\ru{умение} (uméniye) — an ability, know-how}

\label{lesson:RU-C387-umenie}



\begin{quote}
Before the new one: say the Russian for a term, then the Russian for knowledge.
\end{quote}

\subsection*{You'll want to know: \ru{умение}}
\textbf{\ru{умение}} (\emph{uméniye}) — \textquotedblleft{}an ability, know-how\textquotedblright{}.

\begin{quote}
\ru{Умение} \ru{слушать} — \ru{редкий} \ru{дар}.
\end{quote}

\begin{quote}
\emph{Uméniye slúshat — rédkiy dar.}
\end{quote}

\begin{quote}
The ability to listen is a rare gift.
\end{quote}

Neuter. From \emph{\ru{уметь}}, to know how.

\begin{cousinweb}
\textbf{\ru{умение}} is built on \emph{\ru{уметь}} \textquotedblleft{}to know how\textquotedblright{}, which is usually linked with \emph{\ru{ум}} \textquotedblleft{}mind\textquotedblright{}, inherited from Proto-Slavic \emph{*umŭ} (EtymDB).
\end{cousinweb}

\begin{grammarlens}[title={What you've built}]
One more word for telling, explaining and describing at B1.
\end{grammarlens}

\subsection*{Your turn}
\begin{itemize}
\raggedright
  \item \emph{Say it:} \emph{uméniye}
  \item \emph{Say it:} \emph{uméniye}, once more
  \item \emph{Say it:} a sentence of your own with \emph{uméniye}
  \item \emph{Recall:} say the Russian for a term, then the Russian for knowledge, then say \emph{uméniye} again
\end{itemize}

\begin{tcolorbox}[breakable,colback=teal!4,colframe=teal!35!black,title={Before you move on}]
What does \ru{умение} mean? (\textquotedblleft{}An ability, know-how\textquotedblright{}.) Say it once more.
\end{tcolorbox}
\section[\texorpdfstring{\ru{усваивать} (usváivat)}{усваивать (usváivat)}]{\ru{усваивать} (usváivat) — to take in, to master}

\label{lesson:RU-C387-usvaivat}



\begin{quote}
Before the new one: say the Russian for knowledge, then the Russian for an ability.
\end{quote}

\subsection*{You'll want to know: \ru{усваивать}}
\textbf{\ru{усваивать}} (\emph{usváivat}) — \textquotedblleft{}to take in, to master\textquotedblright{}.

\begin{quote}
\ru{Дети} \ru{быстро} \ru{усваивают} \ru{новые} \ru{слова}.
\end{quote}

\begin{quote}
\emph{Déti býstro usváivayut nóvyye slová.}
\end{quote}

\begin{quote}
Children take in new words quickly.
\end{quote}

Imperfective; perfective \emph{\ru{усвоить}}. Built on \emph{\ru{свой}}, one’s own: to make something your own.

\begin{cousinweb}
\textbf{\ru{усваивать}} is the imperfective of \emph{\ru{усвоить}}, built on \emph{\ru{свой}} \textquotedblleft{}one’s own\textquotedblright{}: to master a thing is to make it your own.
\end{cousinweb}

\begin{grammarlens}[title={What you've built}]
One more word for telling, explaining and describing at B1.
\end{grammarlens}

\subsection*{Your turn}
\begin{itemize}
\raggedright
  \item \emph{Say it:} \emph{usváivat}
  \item \emph{Say it:} \emph{usváivat}, once more
  \item \emph{Say it:} a sentence of your own with \emph{usváivat}
  \item \emph{Recall:} say the Russian for knowledge, then the Russian for an ability, then say \emph{usváivat} again
\end{itemize}

\begin{tcolorbox}[breakable,colback=teal!4,colframe=teal!35!black,title={Before you move on}]
What does \ru{усваивать} mean? (\textquotedblleft{}To take in, to master\textquotedblright{}.) Say it once more.
\end{tcolorbox}
\section[\texorpdfstring{\ru{отчислять} (otchislyát)}{отчислять (otchislyát)}]{\ru{отчислять} (otchislyát) — to expel (from a university); to deduct}

\label{lesson:RU-C387-otchislyat}



\begin{quote}
Before the new one: say the Russian for an ability, then the Russian for to take in.
\end{quote}

\subsection*{You'll want to know: \ru{отчислять}}
\textbf{\ru{отчислять}} (\emph{otchislyát}) — \textquotedblleft{}to expel (from a university); to deduct\textquotedblright{}.

\begin{quote}
\ru{Студента} \ru{отчислили} \ru{за} \ru{прогулы}.
\end{quote}

\begin{quote}
\emph{Studénta otchíslili za progúly.}
\end{quote}

\begin{quote}
The student was expelled for missing classes.
\end{quote}

Imperfective; perfective \emph{\ru{отчислить}}. It also means to deduct or transfer money: \emph{\ru{отчислять} \ru{налоги}}.

\begin{cousinweb}
\textbf{\ru{отчислять}} is built from \emph{\ru{от}-} and \emph{\ru{число}} \textquotedblleft{}number\textquotedblright{}, inherited from Proto-Slavic \emph{*čislo} (EtymDB): to take a name off the count.
\end{cousinweb}

\begin{grammarlens}[title={What you've built}]
One more word for telling, explaining and describing at B1.
\end{grammarlens}

\subsection*{Your turn}
\begin{itemize}
\raggedright
  \item \emph{Say it:} \emph{otchislyát}
  \item \emph{Say it:} \emph{otchislyát}, once more
  \item \emph{Say it:} a sentence of your own with \emph{otchislyát}
  \item \emph{Recall:} say the Russian for an ability, then the Russian for to take in, then say \emph{otchislyát} again
\end{itemize}

\begin{tcolorbox}[breakable,colback=teal!4,colframe=teal!35!black,title={Before you move on}]
What does \ru{отчислять} mean? (\textquotedblleft{}To expel (from a university); to deduct\textquotedblright{}.) Say it once more.
\end{tcolorbox}
\section[\texorpdfstring{\ru{образованный} (obrazóvannyy)}{образованный (obrazóvannyy)}]{\ru{образованный} (obrazóvannyy) — educated}

\label{lesson:RU-C387-obrazovannyy}



\begin{quote}
Before the new one: say the Russian for to take in, then the Russian for to expel.
\end{quote}

\subsection*{You'll want to know: \ru{образованный}}
\textbf{\ru{образованный}} (\emph{obrazóvannyy}) — \textquotedblleft{}educated\textquotedblright{}.

\begin{quote}
\ru{Он} \ru{очень} \ru{образованный} \ru{человек}.
\end{quote}

\begin{quote}
\emph{On óchen obrazóvannyy chelovék.}
\end{quote}

\begin{quote}
He is a very educated man.
\end{quote}

Short form \emph{\ru{образован}}. From \emph{\ru{образование}}, education.

\begin{cousinweb}
\textbf{\ru{образованный}} is the participle of \emph{\ru{образовать}} \textquotedblleft{}to form\textquotedblright{}, from \emph{\ru{образ}} \textquotedblleft{}a shape, an image\textquotedblright{}, Proto-Slavic \emph{*obrazŭ} (EtymDB): a person whose mind has been given shape.
\end{cousinweb}

\begin{grammarlens}[title={What you've built}]
One more word for telling, explaining and describing at B1.
\end{grammarlens}

\subsection*{Your turn}
\begin{itemize}
\raggedright
  \item \emph{Say it:} \emph{obrazóvannyy}
  \item \emph{Say it:} \emph{obrazóvannyy}, once more
  \item \emph{Say it:} a sentence of your own with \emph{obrazóvannyy}
  \item \emph{Recall:} say the Russian for to take in, then the Russian for to expel, then say \emph{obrazóvannyy} again
\end{itemize}

\begin{tcolorbox}[breakable,colback=teal!4,colframe=teal!35!black,title={Before you move on}]
What does \ru{образованный} mean? (\textquotedblleft{}Educated\textquotedblright{}.) Say it once more.
\end{tcolorbox}
